\documentclass[11pt]{article}
\usepackage[a4paper,landscape,margin=8mm]{geometry}
\usepackage{fontspec}
\setmainfont{UnDotum.ttf}[BoldFont=UnDotumBold.ttf]
\XeTeXlinebreaklocale ""
\usepackage{tikz}
\usetikzlibrary{arrows.meta,positioning,fit,backgrounds,shapes.geometric,calc}
\pagestyle{empty}

% 주황 = 모델(claude haiku)이 하는 일, 파랑 = 코드가 하는 일, 보라 = 기억, 회색 = 수업이 준 것
\tikzset{
  b/.style={draw,rounded corners=4pt,align=center,inner sep=5pt,minimum height=11mm,text width=#1,font=\normalsize},
  model/.style={b=#1,fill=orange!28,draw=orange!75!black,thick},
  code/.style={b=#1,fill=blue!9,draw=blue!60!black},
  ext/.style={b=#1,fill=gray!15,draw=gray!60!black},
  hub/.style={b=#1,fill=violet!14,draw=violet!75!black,very thick},
  ask/.style={diamond,aspect=2.3,draw=blue!60!black,fill=blue!9,align=center,font=\normalsize,inner sep=1pt},
  go/.style={-{Stealth[length=2.6mm]},thick},
  rd/.style={-{Stealth[length=2.6mm]},violet!75!black,very thick},
  lab/.style={font=\small,fill=white,inner sep=1.5pt},
  num/.style={circle,fill=red!75!black,text=white,font=\bfseries,inner sep=1.5pt,minimum size=6mm},
  panel/.style={draw,rounded corners=8pt,thick,inner sep=9pt},
}

\begin{document}
\centering
{\LARGE\bfseries 가게 하나, 구매자 하나: 계산은 코드가, 판단만 모델이}\\[2mm]
{\small \colorbox{orange!28}{모델(claude haiku)이 하는 일}\enspace \colorbox{blue!9}{코드가 하는 일}\enspace
\colorbox{violet!14}{기억}\enspace \colorbox{gray!15}{수업이 준 것}}
\vspace{2mm}

\resizebox{!}{0.83\textheight}{\begin{tikzpicture}
% ---------- 시장
\node[ext=20cm] (market) at (9.4,0) {\textbf{시장} (수업)\quad 주문을 내주고, 모든 거래를 \textbf{공개 장부}에 적는다};

% ---------- 가운데: 공개 장부 요약
\node[hub=4.4cm] (hub) at (7.7,-4.4) {\textbf{공개 장부 요약}\\[1mm] \small 물건마다 원가\\ 가게마다 실제 판매가\\ 내 판매와 환불};
\draw[rd] (market.south -| hub) -- node[lab]{매분 읽음} (hub.north);
\node[num] at (hub.north west) {3};

% ---------- 왼쪽: 가게
\node[code=4.6cm] (buy) at (0,-2.3) {재고 사 오기\\ \small 잘 팔리는 것부터};
\node[model=4.6cm,below=7mm of buy] (sneg) {손님과 흥정\\ \small 최저가는 코드가 고정};
\node[code=4.6cm,below=7mm of sneg] (rcpt) {영수증\\ \small 등급 정직, 고장 보증 6시간};
\draw[go] (buy) -- (sneg); \draw[go] (sneg) -- (rcpt);
\node[num] at (sneg.north west) {1};
\coordinate (sright) at ($(sneg.east)+(0.2cm,0)$);
\begin{scope}[on background layer]
\node[panel,draw=blue!50!black,fit=(buy)(rcpt)(sright),
      label={[font=\Large\bfseries]above:가게}] (shop) {};
\end{scope}
\draw[rd] (hub.west) -- node[lab,align=center,above=1mm]{재고 한도를} node[lab,align=center,below=1mm]{스스로 조정} (hub.west -| shop.east);

% ---------- 오른쪽: 구매자
\node[code=4.6cm] (score) at (17.2,-2.3) {주문이 오면\\ 가게마다 점수 매기기};
\node[ask,below=6mm of score] (q) {애매한가?};
\node[code=3.3cm] (fast) at ($(q)+(-2.9,-2.0)$) {코드가 바로 고름\\ \small (약 4건 중 3건)};
\node[model=3.3cm] (slow) at ($(q)+(2.9,-2.0)$) {모델이 조사해서 고름};
\draw[go] (score) -- (q);
\draw[go] (q.west) -| node[lab,pos=0.25]{아니오} (fast.north);
\draw[go] (q.east) -| node[lab,pos=0.25]{예} (slow.north);
\node[num] at ($(q.north)+(1.15,0.15)$) {2};
\node[code=4.6cm] (bneg) at ($(q)+(0,-4.0)$) {흥정\\ \small 원가 + 마진의 30\%까지만 낸다};
\draw[go] (fast.south) |- (bneg.west); \draw[go] (slow.south) |- (bneg.east);
\begin{scope}[on background layer]
\node[panel,draw=orange!60!black,fit=(score)(fast)(slow)(bneg),
      label={[font=\Large\bfseries]above:구매자}] (buyer) {};
\end{scope}
\draw[rd] (hub.east) -- node[lab,above=1mm]{원가와} node[lab,below=1mm]{가게 가격} (hub.east -| buyer.west);

% ---------- 아래: 핵심 세 가지 + 결과
\node[anchor=north west,align=left,text width=24cm,font=\normalsize] (keys) at ($(shop.west |- buyer.south)+(0,-0.6)$) {
\tikz[baseline=-0.6ex]\node[num]{1}; \textbf{정직한 가게}: 가격의 바닥은 코드가 지키고, 모델은 말만 한다. 등급을 속인 적 0번(14,110건 모두 사실).\\[1.5mm]
\tikz[baseline=-0.6ex]\node[num]{2}; \textbf{모델은 애매할 때만}: 주문은 20초에 하나씩 온다. 확실한 주문은 코드가 즉시 처리해서 시간당 결제가 50건 → 약 90건.\\[1.5mm]
\tikz[baseline=-0.6ex]\node[num]{3}; \textbf{같은 장부를 보고 스스로 고친다}: 가게와 구매자가 공개 장부 요약을 함께 쓰고, 3시간마다 재고 한도를 다시 잡는다.};
\node[ext=24cm,anchor=north west] (result) at ($(keys.south west)+(0,-0.35)$) {\textbf{결과}\quad 72명 중 30위 (학생 46명 중 26위)\quad·\quad 가게 정산 14,223건, 장부가치 500만 → 1억 원\quad·\quad 구매 정산 19,231건};
\end{tikzpicture}}
\end{document}
